\documentclass{article}

\usepackage[letterpaper, margin=1in]{geometry}
\usepackage{setspace}
\usepackage{amsmath,amssymb}
\usepackage{graphicx}
\usepackage[amsmath,thmmarks,framed]{ntheorem}
\usepackage{authblk}
\usepackage{tcolorbox}
\usepackage{enumerate}
\usepackage{xcolor}
\definecolor{darkgreen}{rgb}{0.0,0.5,0.0}
\usepackage[colorlinks,linkcolor=red,citecolor=darkgreen,urlcolor=blue]{hyperref}
\usepackage[round,authoryear]{natbib}

\usepackage[ruled,linesnumbered]{algorithm2e}
\SetKwInput{KwInput}{Input}
\SetKwInput{KwOutput}{Output}

\newtheorem*{remark}{Remark}
\newtheorem{lemma}{Lemma}
\newenvironment{proof}{\par\medskip\noindent\textit{Proof.}\ }{\hfill$\square$\par\medskip}

\newcommand{\Le}{\leqslant}
\newcommand{\Ge}{\geqslant}
\newcommand{\D}{\mathrm{d}}
\newcommand{\DB}{\mathrm{DB}}
\newcommand{\KL}{\mathrm{KL}}
\newcommand{\TV}{\mathrm{TV}}
\newcommand{\X}{\mathrm{X}}
\newcommand{\ESS}{\mathrm{ESS}}
\newcommand{\setzero}{\setlength{\itemsep}{0pt}}

\theoremstyle{plain}

\title{Beyond Fake Effective Sample Size: Normalizing Flow Boltzmann Sampling with Initiative Mode Discovery}
\author[1]{Qi Feng}
\author[2]{Rongjie Lai}
\author[2]{Di Qi}
\author[2]{Xuda Ye}
\affil[1]{Department of Mathematics, Florida State University, Tallahassee, FL 32306}
\affil[2]{Department of Mathematics, Purdue University, West Lafayette, IN 47907}

\date{\today}

\begin{document}
\maketitle

\section{Introduction}

The Boltzmann generator built on normalizing flows \citep{noe2019boltzmann,wu2020stochastic} is a central tool for sampling from complicated distributions. Let $\mu$ denote the target distribution on $\mathbb R^d$ and let $\nu$ denote the pushforward distribution of the flow. The conventional Boltzmann generator trains the flow by minimizing the reverse KL
\begin{equation}
    \mathrm{KL}(\nu\|\mu) = \int_{\mathbb R^d} \nu(x) \log \frac{\nu(x)}{\mu(x)} \D x
    \label{eq: reverse KL}
\end{equation}
with respect to the flow parameters.
On multimodal targets, however, flows trained with the reverse KL are prone to mode collapse: any mode lost during training cannot be recovered, severely degrading sample quality. A standard remedy is to replace the reverse KL with the forward KL \citep{gabrie2022adaptive}:
\begin{equation}
    \mathrm{KL}(\mu\|\nu) = \int_{\mathbb R^d} \mu(x) \log \frac{\mu(x)}{\nu(x)} \D x,
    \label{eq: forward KL}
\end{equation}
which is more on multimodal targets. The forward KL is the expectation of the \emph{log-ratio} $\log(\mu/\nu)$ under $\mu$, and therefore requires samples from the target. Such samples can be approximated by annealed importance sampling (AIS) \citep{midgley2022flow}, which transports the pushforward $\nu$ toward $\mu$ through a sequence of reweighting and resampling steps.
The flow AIS bootstrap (FAB) \citep{midgley2022flow} follows the same idea but trains against the more mode-sensitive weight $\mu^2/\nu$ in place of $\mu$, enhancing mode discovery.

These objectives can in principle mitigate mode collapse, but their effectiveness hinges on AIS samples that faithfully cover the target. If the pushforward $\nu$ misses modes of $\mu$, the AIS chain---initialized at $\nu$---is unlikely to reach them, and the resulting samples remain biased toward the modes $\nu$ already covers. This exposes a fatal defect of any AIS-reliant method, including the forward KL and FAB: the effective sample size (ESS) \citep{doucet2001introduction}, the gold-standard metric for assessing flow quality, can be \emph{fake}. An ESS close to $100\%$ does \emph{not} guarantee mode coverage, since modes missing from $\nu$ contribute no mass to the ESS calculation, and their absence goes undetected.

This paper targets the fake ESS problem and develops a training scheme with \emph{initiative mode discovery}, in which prior knowledge of the modes of $\mu$ is injected directly into the loss to guide the flow. Doing this safely is far from trivial: the forward KL and FAB are unbiased only when fed \emph{accurate} samples from $\mu$, so naively augmenting the AIS samples with the known mode locations is highly likely to corrupt the minimizer---the injected samples carry no correct importance weights and therefore pull the optimum away from $\mu$.

The center of our training scheme is the cross-regularization functional (denoted $\X$),
\begin{equation}
    \X_\omega(\mu\|\nu) = \iint_{\mathbb R^d\times\mathbb R^d} \omega(x,y)
    \biggl|\log\frac{\mu(x)}{\nu(x)} - \log\frac{\mu(y)}{\nu(y)}\biggr| \D x \D y,
    \label{eq: X}
\end{equation}
where $\omega(x,y)$ is a finite positive measure on $\mathbb R^{2d}$ whose choice we defer to a later section. The motivation is direct: since our goal is to match $\nu$ with $\mu$, the natural quantity to minimize is the log-ratio $\log(\mu/\nu)$ itself. We penalize the \emph{difference} of two log-ratios rather than a single absolute value because $\mu$ is known only up to an unknown normalization constant, and that constant cancels exactly in the difference. $\X$ offers three key benefits as a loss-function ingredient:
\begin{enumerate}
    \setzero
    \item \textbf{Pointwise penalty.} The forward KL reshapes $\nu$ globally through Jensen's inequality, whereas $\X$ penalizes log-ratio variations wherever they occur, making it sensitive to local features of the potential landscape that purely global divergences may miss.

    \item \textbf{Tolerance to inaccurate target samples.} $\X$ vanishes whenever the log-ratio is pointwise constant, regardless of the accuracy of the samples drawn from $\mu$. Inaccuracies in the AIS samples therefore barely shift the minimizer, and manually injected mode locations stay safe inside the loss---unlocking the initiative mode discovery promised above.

    \item \textbf{Drop-in compatibility with the forward KL.} $\X$ shares the integrand $\log\nu$ with the forward KL, so when $\omega(x,y) = \mu(x)\mu(y)$ the same autograd graph and the same $\mu$-samples are reused at negligible extra cost; moreover $\X$ has the same scale as the forward KL, so no additional tuning is required.
\end{enumerate}

We consider two choices of the measure $\omega(x,y)$ in the X functional \eqref{eq: X}, each playing a distinct role. The first is $\omega(x,y) = \mu(x)\mu(y)$, giving
\begin{equation}
    \X_\mu(\mu\|\nu) = \iint_{\mathbb R^d\times\mathbb R^d}
    \mu(x) \mu(y) \biggl|\log\frac{\mu(x)}{\nu(x)} - \log\frac{\mu(y)}{\nu(y)}\biggr| \D x \D y.
    \label{eq: X-mu}
\end{equation}
As noted above, computing $\X_\mu$ and its gradient incurs almost no extra cost beyond the forward KL, so $\X_\mu$ can be added as a stability- and accuracy-improving regularizer at essentially zero overhead.

The second choice is $\omega(x,y) = \hat\mu(x)\hat\mu(y)$, sharing the same form as \eqref{eq: X-mu}:
\begin{equation}
    \X_{\hat\mu}(\mu\|\nu) = \iint_{\mathbb R^d\times\mathbb R^d}
    \hat\mu(x) \hat\mu(y) \biggl|\log\frac{\mu(x)}{\nu(x)} - \log\frac{\mu(y)}{\nu(y)}\biggr| \D x \D y,
    \label{eq: X-hat-mu}
\end{equation}
where $\hat\mu$ \emph{need not be an accurate approximation to $\mu$}: by the second benefit of $\X$, the minimizer of $\X_{\hat\mu}$ is insensitive to the precise shape of $\hat\mu$. The sole role of $\hat\mu$ is to place mass on modes that the current pushforward $\nu$ has not yet reached, so a coarse mode-finder that errs on the side of broad coverage is exactly what we want. We construct such an $\hat\mu$ via the \emph{Quench and Temper} (QT) algorithm, executed in three steps:
\begin{enumerate}
    \setzero
    \item \emph{Diffusion} (melt): add large Brownian noise to the samples so that they cover a wider region of $\mathbb R^d$;
    \item \emph{Optimization} (quench): apply LBFGS to drive the noisy samples toward the mode centers of $\mu$;
    \item \emph{Rejuvenation} (temper): apply a short-time Langevin so that the samples relax across the local landscape of each mode.
\end{enumerate}
The three steps are complementary: the Brownian noise enables discovery of modes not yet covered by $\nu$; LBFGS locates the mode centers precisely; the Langevin step diffuses the samples and makes $\hat\mu$ robust to small perturbations. The name borrows from the metallurgical analogue: quenching rapidly cools the material to lock in a low-energy configuration, while tempering gently reheats it to relieve stress and reach a more stable state. The procedure also applies to targets whose mode set forms a low-dimensional manifold, since a finite-time optimization can still capture such a manifold.

Combining the forward KL with both regularizers yields our training loss for initiative mode discovery:
\begin{align}
    \mathcal L[\nu] & = \KL(\mu\|\nu) +
    \lambda\, \X_\mu(\mu\|\nu) + \hat\lambda\, \X_{\hat\mu}(\mu\|\nu) \nonumber \\
    & = \int_{\mathbb R^d} \mu(x) \log \frac{\mu(x)}{\nu(x)} \D x +
    \iint_{\mathbb R^d\times\mathbb R^d}
    \Bigl(\lambda\, \mu(x)\mu(y) + \hat\lambda\, \hat\mu(x)\hat\mu(y)\Bigr) \biggl|\log\frac{\mu(x)}{\nu(x)} - \log\frac{\mu(y)}{\nu(y)}\biggr| \D x \D y,
\end{align}
where $\lambda, \hat\lambda \Ge 0$ are tunable weights. To evaluate the trained flow we report ESS \emph{together with} the coverage metric \citep{naeem2020reliable}, which measures how many clusters of $\hat\mu$ are reached by samples drawn through $\nu$. The ESS alone can be fake when modes of $\mu$ are missing, so it must be paired with a mode-aware diagnostic: only when both ESS and coverage are high does the flow truly approximate $\mu$.

\section{Fisher--Rao gradient flow with X functionals}
\label{sec: FR}

Before we propose the training details with the $\X$ functionals, we first study the gradient flow properties of the loss function, whose convergence rate reflects the difficulty with which the flow can be trained. We view $\nu$ as a generic probability density and ask how the loss $\mathcal L[\nu]$ decreases under the gradient flow induced by an appropriate Riemannian metric on the space of densities.

Consider the loss function with respect to $\nu$ with only the forward KL and $\X_\mu$:
\begin{equation}
    \mathcal L[\nu] = \KL(\mu\|\nu) + \lambda\, \X_\mu(\mu\|\nu) = \int_{\mathbb R^d} \mu(x) \log \frac{\mu(x)}{\nu(x)} \D x +
    \iint_{\mathbb R^d\times\mathbb R^d}
    \lambda\, \mu(x)\mu(y) \biggl|\log\frac{\mu(x)}{\nu(x)} - \log\frac{\mu(y)}{\nu(y)}\biggr| \D x \D y.
    \label{eq: loss KL X}
\end{equation}
Throughout this section we abbreviate the importance weight by
\begin{equation*}
    w(x) := \frac{\mu(x)}{\nu(x)},
\end{equation*}
so that the log-ratio is $\log w(x) = \log\mu(x) - \log\nu(x)$.

The choice of Riemannian metric matters. In the Wasserstein-2 geometry, the forward KL fails to be displacement convex along W$_2$ geodesics, and the standard Otto--Villani / Bakry--{\'E}mery machinery that yields exponential convergence for the \emph{reverse} KL (Langevin / Fokker--Planck) does not apply. We therefore work in the \emph{Fisher--Rao} (FR) geometry, in which the forward KL admits an unconditional exponential decay rate and the regularizer $\X_\mu$ provides an additional, strictly positive boost to that rate.

\subsection{First variations}

The first variations of the two summands of $\mathcal L$ with respect to $\nu$ admit a clean closed form. For the forward KL,
\begin{equation}
    \frac{\delta\,\KL(\mu\|\nu)}{\delta \nu}(z) = -\frac{\mu(z)}{\nu(z)}.
    \label{eq: var-KL}
\end{equation}

For the regularizer, note that $\delta \log w(x)/\delta\nu(z) = -\delta(x-z)/\nu(x)$. Differentiating $|\log w(x) - \log w(y)|$ with the subdifferential $\partial|s|/\partial s = \operatorname{sgn}(s)$, using $\operatorname{sgn}(\log w(x) - \log w(y)) = \operatorname{sgn}(w(x) - w(y))$, and exploiting the antisymmetry of $\operatorname{sgn}$ to combine the two resulting terms gives
\begin{equation}
    \frac{\delta\,\X_\mu(\mu\|\nu)}{\delta \nu}(z) = -\frac{2\mu(z)}{\nu(z)}\,S(z),
    \qquad
    S(z) := \mathbb E_{Y\sim\mu}\bigl[\operatorname{sgn}(w(z) - w(Y))\bigr]\in [-1,1].
    \label{eq: var-X}
\end{equation}
We call $S$ the \emph{rank field} of the importance weight under $\mu$: $S(z) > 0$ where $w(z)$ exceeds the $\mu$-median of $w$, i.e. where $\nu$ undercovers a mode of $\mu$; $S(z) < 0$ where $\nu$ overcovers; and $S \equiv 0$ at $\nu = \mu$. By the antisymmetry $\operatorname{sgn}(w(X) - w(Y)) = -\operatorname{sgn}(w(Y) - w(X))$ under $X\leftrightarrow Y$ swap with $X, Y\sim\mu$, we also have $\mathbb E_\mu[S] = 0$.

Combining \eqref{eq: var-KL} and \eqref{eq: var-X},
\begin{equation*}
    \frac{\delta \mathcal L}{\delta\nu}(z) = -\frac{\mu(z)}{\nu(z)}\bigl(1 + 2\lambda S(z)\bigr),
\end{equation*}
so the factor $1 + 2\lambda S \in [1-2\lambda, 1+2\lambda]$ acts as a rank-based modulator of the pure forward-KL drift.

\subsection{The Fisher--Rao gradient flow}

The Fisher--Rao metric on the space of probability densities is the Riemannian metric
\begin{equation*}
    \langle \sigma_1, \sigma_2\rangle^{\mathrm{FR}}_\nu = \int_{\mathbb R^d} \frac{\sigma_1(x)\,\sigma_2(x)}{\nu(x)}\,\D x,
\end{equation*}
where tangent vectors $\sigma_i$ are signed measures with $\int\sigma_i = 0$. Under the change of variable $\nu\mapsto 2\sqrt{\nu}$, this metric pulls back to the standard $L^2$ metric on a piece of the unit sphere; FR geodesics interpolate $\sqrt{\nu}$ linearly (with renormalization), equivalently they interpolate $\log\nu$ affinely modulo a constant. The geometry is \emph{purely multiplicative}: mass at $x$ is reweighted, never transported.

For a functional $\mathcal F[\nu]$ with first variation $\delta\mathcal F/\delta\nu$, the FR gradient is defined by the identity
\begin{equation*}
    \langle\mathrm{grad}^{\mathrm{FR}}\mathcal F,\,\sigma\rangle^{\mathrm{FR}}_\nu = \int_{\mathbb R^d} \frac{\delta\mathcal F}{\delta\nu}(x)\,\sigma(x)\,\D x \quad \text{for all admissible } \sigma.
\end{equation*}
Writing $\mathrm{grad}^{\mathrm{FR}}\mathcal F(x) = \nu(x)\,h(x)$, this forces $h = \delta\mathcal F/\delta\nu - c$, and the constraint $\int\mathrm{grad}^{\mathrm{FR}}\mathcal F = 0$ fixes the constant $c = \langle\delta\mathcal F/\delta\nu\rangle_\nu$, where we denote the $\nu$-mean by
\begin{equation*}
    \langle\,\cdot\,\rangle_\nu := \int_{\mathbb R^d}(\,\cdot\,)\,\nu(x)\,\D x.
\end{equation*}
Hence
\begin{equation}
    \mathrm{grad}^{\mathrm{FR}}\mathcal F(\nu) = \nu\Bigl(\frac{\delta\mathcal F}{\delta\nu} - \big\langle\frac{\delta\mathcal F}{\delta\nu}\big\rangle_\nu\Bigr).
    \label{eq: FR-gradient}
\end{equation}
The FR gradient flow $\partial_t\nu_t = -\mathrm{grad}^{\mathrm{FR}}\mathcal F$ automatically conserves probability mass and admits the universal dissipation identity
\begin{equation*}
    \frac{\D}{\D t}\mathcal F[\nu_t] = -\bigl\|\mathrm{grad}^{\mathrm{FR}}\mathcal F\bigr\|^2_{\mathrm{FR},\nu_t} = -\!\int\Bigl(\frac{\delta\mathcal F}{\delta\nu} - \big\langle\frac{\delta\mathcal F}{\delta\nu}\big\rangle_{\nu_t}\Bigr)^{\!2}\nu_t\,\D x.
\end{equation*}

\subsection{Exponential convergence of the forward KL flow}
\label{subsec: FR-KL}

Applying \eqref{eq: FR-gradient} to the forward KL with $\delta\KL/\delta\nu = -\mu/\nu$ and $\langle -\mu/\nu\rangle_\nu = -\!\int\mu = -1$ gives the strikingly simple flow
\begin{equation}
    \partial_t\nu_t = -\mathrm{grad}^{\mathrm{FR}}\KL(\nu_t) = \mu - \nu_t,
    \label{eq: FR-KL-flow}
\end{equation}
a linear ODE on densities that decouples pointwise in $x$. Its closed-form solution is
\begin{equation*}
    \nu_t(x) = e^{-t}\,\nu_0(x) + (1-e^{-t})\,\mu(x).
\end{equation*}

Differentiating the KL along this flow,
\begin{equation}
    \frac{\D}{\D t}\KL(\mu\|\nu_t) = -\!\int_{\mathbb R^d}\!\frac{\mu}{\nu_t}\,(\mu - \nu_t)\,\D x = 1 - \!\int_{\mathbb R^d}\!\frac{\mu^2}{\nu_t}\,\D x = -\,\chi^2(\mu\|\nu_t),
    \label{eq: KL-dissipation-pure}
\end{equation}
where the chi-squared divergence is
\begin{equation*}
    \chi^2(\mu\|\nu) := \int_{\mathbb R^d}(w - 1)^2\,\nu\,\D x = \int_{\mathbb R^d} w^2\,\nu\,\D x - 1.
\end{equation*}
The elementary inequality $u\log u \Le u^2 - u$ applied pointwise to $u = w$ yields $\chi^2(\mu\|\nu) \Ge \KL(\mu\|\nu)$, and Gr\"onwall's lemma then gives
\begin{equation}
    \KL(\mu\|\nu_t) \Le e^{-t}\,\KL(\mu\|\nu_0),\qquad t \Ge 0.
    \label{eq: KL-exponential}
\end{equation}

\begin{remark}
The exponential rate $e^{-t}$ holds with \emph{no structural assumption on the target $\mu$} --- neither log-concavity, nor a spectral gap, nor any growth condition. This contrasts sharply with the W$_2$ geometry, where the forward KL is not displacement convex and an analogous exponential rate is unavailable without further assumptions.
\end{remark}

The uniform convergence rate \eqref{eq: KL-exponential} should \emph{not} be read as a blanket statement that minimizing the forward KL is intrinsically faster or more efficient than minimizing the reverse KL. The FR flow is a continuous-time idealization on the space of densities and is agnostic to the sampling cost. In practice, the difficulty of training has been partially transported to the AIS step: evaluating the forward KL and its gradient requires accurate samples of $\mu$, and producing such samples in high dimensions is itself a hard problem --- it is the very reason that AIS is needed, and that the $\X_{\hat\mu}$ regularizer must compensate for modes that AIS alone fails to reach. The exponential rate \eqref{eq: KL-exponential} quantifies how fast the flow \emph{can} converge given accurate target samples; it does not eliminate the burden of producing them.

\subsection{Acceleration via $\X_\mu$}
\label{subsec: FR-X-acceleration}

We now compute the FR gradient of the regularizer $\X_\mu$. Using \eqref{eq: var-X} and the symmetry $\mathbb E_\mu[S] = 0$,
\begin{equation*}
    \langle -2(\mu/\nu)\,S\rangle_\nu = -2\!\int\mu\,S\,\D z = -2\,\mathbb E_\mu[S] = 0,
\end{equation*}
so the mean-subtraction in \eqref{eq: FR-gradient} acts trivially, and
\begin{equation}
    \mathrm{grad}^{\mathrm{FR}}\X_\mu(\nu) = -2\mu\,S,
    \qquad
    \partial_t\nu_t = -\mathrm{grad}^{\mathrm{FR}}\X_\mu(\nu_t) = 2\mu\,S(\nu_t).
    \label{eq: FR-X-flow}
\end{equation}
The flux is a fixed multiple of $\mu$ modulated by the rank field $S$: mass is added where $\nu$ undercovers $\mu$ and removed where $\nu$ overcovers $\mu$. Probability is conserved since
\begin{equation*}
    \int_{\mathbb R^d} 2\mu(x)\,S(x)\,\D x = 2\,\mathbb E_\mu[S] = 0.
\end{equation*}

Combining \eqref{eq: FR-KL-flow} and \eqref{eq: FR-X-flow}, the FR gradient flow of the full loss \eqref{eq: loss KL X} is
\begin{equation}
    \partial_t\nu_t = \mu\bigl(1 + 2\lambda\, S(\nu_t)\bigr) - \nu_t.
    \label{eq: FR-L-flow}
\end{equation}
Mass conservation follows from $\int[\mu(1+2\lambda S) - \nu_t]\,\D x = (1 + 2\lambda\mathbb E_\mu[S]) - 1 = 0$, and the flow is stationary at $\nu_t = \mu$, where $w\equiv 1$ and hence $S\equiv 0$.

The acceleration provided by $\X_\mu$ is captured by the KL dissipation along the combined flow:
\begin{equation}
    \frac{\D}{\D t}\KL(\mu\|\nu_t) = -\!\int\frac{\mu}{\nu_t}\bigl[\mu(1 + 2\lambda S) - \nu_t\bigr]\D x = -\chi^2(\mu\|\nu_t) - 2\lambda\!\int_{\mathbb R^d}\frac{\mu^2}{\nu_t}\,S\,\D x.
    \label{eq: KL-dissipation-combined}
\end{equation}
The first term reproduces the pure-KL dissipation \eqref{eq: KL-dissipation-pure}; the second is the contribution of $\X_\mu$, whose sign is settled by the following identity.

\begin{lemma}
\label{lem: rank-correlation}
For any probability density $\nu$ on $\mathbb R^d$,
\begin{equation}
    \int_{\mathbb R^d}\frac{\mu(z)^2}{\nu(z)}\,S(z)\,\D z = \frac{1}{2}\,\mathbb E_{X,Y\stackrel{\mathrm{iid}}{\sim}\mu}\bigl[\,|w(X) - w(Y)|\,\bigr] \;\Ge\; 0,
\end{equation}
with equality if and only if $w$ is $\mu$-almost-everywhere constant, i.e. $\nu = \mu$.
\end{lemma}

\begin{proof}
By the definition of $S$,
\begin{equation*}
    \int_{\mathbb R^d}\frac{\mu^2}{\nu}\,S\,\D z = \mathbb E_\mu\bigl[w(X)\,S(X)\bigr] = \mathbb E_{X,Y\sim\mu}\bigl[w(X)\operatorname{sgn}(w(X) - w(Y))\bigr].
\end{equation*}
Symmetrizing in $X\leftrightarrow Y$ and averaging the two equal copies of the expectation gives
\begin{equation*}
    \int_{\mathbb R^d}\frac{\mu^2}{\nu}\,S\,\D z = \frac{1}{2}\,\mathbb E_{X,Y\sim\mu}\bigl[(w(X) - w(Y))\operatorname{sgn}(w(X) - w(Y))\bigr] = \frac{1}{2}\,\mathbb E_{X,Y\sim\mu}\bigl[\,|w(X) - w(Y)|\,\bigr].
\end{equation*}
\end{proof}

Writing $w_t := \mu/\nu_t$ and combining Lemma~\ref{lem: rank-correlation} with \eqref{eq: KL-dissipation-combined} and $\chi^2 \Ge \KL$,
\begin{equation*}
    \frac{\D}{\D t}\KL(\mu\|\nu_t) \Le -\chi^2(\mu\|\nu_t) - \lambda\,\mathbb E_{X,Y\sim\mu}\bigl[\,|w_t(X) - w_t(Y)|\,\bigr] \Le -\KL(\mu\|\nu_t).
\end{equation*}
The bound \eqref{eq: KL-exponential} therefore \emph{continues to hold along the combined flow}, and the rate strictly improves whenever the importance weights $w_t$ are non-constant under $\mu$, i.e. whenever $\nu_t \ne \mu$. The additional dissipation contributed by $\X_\mu$,
\begin{equation*}
    \Delta_\lambda(\nu_t) := \lambda\,\mathbb E_{X,Y\sim\mu}\bigl[\,|w_t(X) - w_t(Y)|\,\bigr] \;\Ge\; 0,
\end{equation*}
is the precise rate of acceleration: it is strictly positive away from equilibrium, scales linearly in $\lambda$, and quantifies the mean absolute spread of importance weights under $\mu$. The regularizer acts most strongly when the weights are highly variable across the support of $\mu$ --- exactly the regime in which the pushforward $\nu_t$ unevenly covers the modes of $\mu$ --- and it vanishes precisely at the equilibrium $\nu_t = \mu$. Thus adding $\X_\mu$ to the forward KL accelerates the Fisher--Rao gradient flow without altering the unconditional exponential rate guaranteed by the pure KL flow.

\begin{remark}
The regularizer $\X_\mu$ also decreases monotonically along the combined flow: a direct computation along \eqref{eq: FR-L-flow} gives
\begin{equation*}
    \frac{\D}{\D t}\X_\mu(\mu\|\nu_t) = -2\!\int_{\mathbb R^d}\frac{\mu^2}{\nu_t}\,S\,\D x - 4\lambda\!\int_{\mathbb R^d}\frac{\mu^2}{\nu_t}\,S^2\,\D x \;\Le\; 0,
\end{equation*}
with strict inequality whenever $\nu_t \ne \mu$. The regularizer therefore not only accelerates KL convergence but is itself driven to zero in the long-time limit.
\end{remark}

\section{Flow training with X functionals}

In this section we derive the total training loss that drives the normalizing flow from the source distribution to the target, combining the forward KL with the $\X$ regularizers \eqref{eq: X-mu} and \eqref{eq: X-hat-mu}. Let $U_0$ and $U_1$ be the source and target potential functions on $\mathbb R^d$, and let $\mu_0 \propto \exp(-U_0)$ and $\mu_1 \propto \exp(-U_1)$ be the corresponding distributions. Given sufficient samples from $\mu_0$, we seek a differentiable flow map $G : \mathbb R^d \to \mathbb R^d$ such that $G_{\#}\mu_1 \approx \mu_0$; target samples are then drawn from $G^{-1}_{\#}\mu_0$ via importance sampling.

The remainder of this section is organized as follows. We first review the forward KL loss and describe how annealed importance sampling (AIS) supplies approximate samples from the target $\mu_1$. We then detail the computation of $\X_\mu$ \eqref{eq: X-mu}, including the autograd-reuse trick noted earlier. Next, we present the Quench and Temper algorithm that constructs the wide-coverage distribution $\hat\mu$ used in $\X_{\hat\mu}$ \eqref{eq: X-hat-mu}. Finally, we define the ESS and coverage metrics and use a 2D example to demonstrate the power of the $\X$-augmented loss.

\subsection{Forward KL and annealed importance sampling}

By the change-of-variables formula, the pullback density of $\mu_0$ under $G$ is
\begin{equation}
    (G^{-1}_{\#}\mu_0)(y) = \mu_0(G(y)) \, |\det J_{G}(y)|,
    \qquad y \in \mathbb R^d,
\end{equation}
where $J_G(y) \in \mathbb R^{d\times d}$ is the Jacobian of $G$ at $y$. The forward KL loss then reads
\begin{align}
    \KL\big(\mu_1\|G^{-1}_{\#}\mu_0\big)
    & = -\int_{\mathbb R^d} \mu_1(y) \log (G^{-1}_{\#}\mu_0)(y) \D y + \mathrm{const} \notag \\
    & = \int_{\mathbb R^d} \mu_1(y) \Big(U_0(G(y)) - U_1(y) - \log|\det J_G(y)|\Big) \D y + \mathrm{const}.
    \label{eq: forward KL full}
\end{align}
We deliberately parameterize the loss in terms of $G$ (mapping target to source) rather than $G^{-1}$, so that no inversion is required during training.

Evaluating \eqref{eq: forward KL} requires samples from $\mu_1$, which are not directly available. The natural surrogate is to draw $y_1,\dots,y_B$ from the pushforward $G^{-1}_{\#}\mu_0$ and reweight by
\begin{equation}
    w(y) = \frac{\mu_1(y)}{(G^{-1}_{\#}\mu_0)(y)} = \exp\Big(
        U_0(G(y)) - U_1(y) - \log|\det J_G(y)|
    \Big), \qquad y \in \mathbb R^d,
    \label{eq: weight}
\end{equation}
followed by a resampling step and a short Langevin simulation to diffuse the duplicates. In high dimensions, however, $\mu_1$ and $G^{-1}_{\#}\mu_0$ typically differ greatly, so the weights $w(y_i)$ collapse onto a few outliers and the resulting samples are of very low quality.

Annealed importance sampling (AIS) bridges this gap by inserting $M+1$ intermediate distributions
\begin{equation}
    \pi_k = (G^{-1}_{\#}\mu_0)^{1-\frac{k}{M}} (\mu_1)^{\frac{k}{M}} = (G^{-1}_{\#}\mu_0) \cdot w(y)^{\frac{k}{M}},
    \qquad k = 0,1,\dots,M,
    \label{eq: pi-k}
\end{equation}
so that $\pi_0 = G^{-1}_{\#}\mu_0$ and $\pi_M = \mu_1$. Each rung of the ladder reweights by only $w(y)^{\frac{1}{M}}$, keeping the weight ratio close to $1$ and avoiding the degeneracy of direct importance sampling:

\begin{algorithm}[H]
\setstretch{1.2}
\caption{Annealed importance sampling from $G_{\#}^{-1}\mu_0$ to $\mu_1$}
\KwInput{samples $y_1,\dots,y_B$ from $G_{\#}^{-1}\mu_0$, ladder length $M$}
\KwOutput{samples $y_1,\dots,y_B$ approximately distributed as $\mu_1$}

\For{$k = 1,\dots,M$}{
    Compute the per-step weights $w(y_i)^{\frac{1}{M}}$ on $y_1,\dots,y_B$\;
    Resample $y_1',\dots,y_B'$ from the weighted empirical distribution\;
    \tcp{The resampled set has duplicates; Langevin diffuses them}
    Apply a short Langevin simulation to $y_1',\dots,y_B'$, returning fresh samples $y_1,\dots,y_B \sim \pi_k$\;
}
\end{algorithm}

\subsection{Single-pass computation of forward KL and $\X_\mu$}

Computing $\X_\mu$ reuses almost all of the forward KL machinery. Substituting the explicit log-ratio under the flow map $G$ (the subscript $\mu$ here refers to the target $\mu_1$), the $\X_\mu$ functional \eqref{eq: X-mu} becomes
\begin{align}
    \X_\mu\big(\mu_1 \| G^{-1}_{\#}\mu_0\big)
    &= \iint_{\mathbb R^d\times\mathbb R^d}
    \mu_1(y) \mu_1(y') \biggl|
    \Big(U_0(G(y)) - U_1(y) - \log|\det J_G(y)|\Big) \notag \\
    &\qquad{} - \Big(U_0(G(y')) - U_1(y') - \log|\det J_G(y')|\Big)
    \biggr| \D y \D y'.
    \label{eq: X-mu full}
\end{align}
Both the forward KL \eqref{eq: forward KL} and $\X_\mu$ \eqref{eq: X-mu full} can therefore be evaluated in a single forward pass that shares the same autograd graph:

\begin{algorithm}[H]
\setstretch{1.2}
\caption{Single-pass computation of forward KL and $\X_\mu$}
\KwInput{batch $\mathbf y = (y_1,\dots,y_B) \in \mathbb R^{B\times d}$ of samples from $\mu_1$}
\KwOutput{forward KL term and $\X_\mu$ term}
\tcp{$\mathbf z$ carries the autograd graph through $G$}
Compute $\mathbf z = U_0(G(\mathbf y)) - U_1(\mathbf y) - \log|\det J_G(\mathbf y)|$\;
\tcp{$\mathbf z'$ inherits the graph of $\mathbf z$ at zero cost}
Sample a random permutation $\sigma$ of $\{1,\dots,B\}$ and set $\mathbf z' = \mathbf z[\sigma]$\;
Return $\mathrm{mean}(\mathbf z)$ as the forward KL term and $\mathrm{mean}\bigl(|\mathbf z - \mathbf z'|\bigr)$ as the $\X_\mu$ term\;
\end{algorithm}
\mbox{}\\
Consequently, the only autograd work touching the flow parameters is the single backward pass through $\mathbf z$, identical to a vanilla forward KL step. Adding $\X_\mu$ therefore incurs essentially no additional training cost.

\subsection{Initiative mode discovery in $X_{\hat \mu}$ with QT algorithm}

Computing the $\X_{\hat\mu}$ functional \eqref{eq: X-hat-mu} requires samples from $\hat\mu_1$, the wide-coverage approximation of the target $\mu_1$. These samples are prepared once, before training: we generate a fixed pool of $P$ samples $x_1,\dots,x_P$ distributed as $\hat\mu_1$, and at each training step we draw a mini-batch of size $B$ from the pool and refresh it with a short Langevin chain before evaluating the loss. The pool is produced by the QT algorithm, which proceeds in three stages: \emph{diffusion} (melt), \emph{optimization} (quench), and \emph{rejuvenation} (temper).

\begin{algorithm}[H]
\label{alg: QT}
\setstretch{1.2}
\caption{QT: initiative mode discovery}
\KwInput{samples $x_1,\dots,x_P$ from $\mu_0$}
\KwOutput{samples $x_1,\dots,x_P$ forming $\hat\mu_1$}
\tcp{Diffusion (melt): scatter the samples to a wider region of $\mathbb R^d$}
Choose a large $\sigma > 0$ and update $x_i \leftarrow x_i + \sigma\xi_i$ with $\xi_i \sim \mathcal N(0, I_d)$ i.i.d.\;
\tcp{Optimization (quench): drive each sample toward a mode center of $\mu_1$}
Run LBFGS on $U_1$ for a fixed number of steps to update $x_1,\dots,x_P$\;
\tcp{Rejuvenation (temper): spread the samples around each mode}
Run a short Langevin simulation targeting $\mu_1$ on $x_1,\dots,x_P$\;
\end{algorithm}
\mbox{}\\
The output samples concentrate near the mode centers of $\mu_1$, and the initial Brownian diffusion ensures that even modes far from the support of the input are reachable. The $\X_{\hat\mu}$ regularizer therefore penalizes any mode of $\mu_1$ on which the current pushforward $G^{-1}_{\#}\mu_0$ places too little mass, driving the flow to discover modes that AIS alone could never reach.

A side effect of the diffusion step is that a large $\sigma$ biases the output samples toward the outer modes---those far from the bulk of $\mu_0$. This bias is not only harmless but actually desirable for mode discovery, since the modes that the current pushforward fails to cover are precisely those lying far from $\mu_0$.

\subsection{Evaluation metrics: ESS and coverage}

The trained flow $G$ induces the distribution $G^{-1}_{\#}\mu_0$ as an approximation to the target $\mu_1$. We assess its quality through two complementary metrics: the effective sample size (ESS), which measures how uniform the importance weights are on the support reached by the flow, and the coverage, which checks whether every mode of $\mu_1$ has been reached at all.

The first metric is the effective sample size. Let $\mathbf y = (y_1,\dots,y_N) \in \mathbb R^{N\times d}$ be i.i.d.\ samples from $G^{-1}_{\#}\mu_0$, with importance weights $w(y_i)$ given by \eqref{eq: weight}. The ESS is defined as
\begin{equation}
    \ESS(\mathbf y) = \frac{\big(\sum_{i=1}^N w(y_i)\big)^2}{N \sum_{i=1}^N w(y_i)^2} \in [0,1],
\end{equation}
where the upper bound $1$ follows from the Cauchy--Schwarz inequality and is attained when all weights are equal. When $G$ is perfectly trained the weights are constant and the ESS reaches $1$. As emphasized in the introduction, however, this conclusion is only one-way: a high ESS certifies that the weights are uniform \emph{on the support of $G^{-1}_{\#}\mu_0$}, but modes of $\mu_1$ that the pushforward never reaches contribute nothing to the sum and so cannot lower the ESS---hence the \emph{fake ESS} pitfall.

To detect such missing modes, we pair the ESS with the coverage metric of \citet{naeem2020reliable}. Let $\mathbf x = (x_1,\dots,x_P) \in \mathbb R^{P\times d}$ be the samples of $\hat\mu_1$ produced by the Quench and Temper algorithm (Algorithm~\ref{alg: QT}). For each $1 \Le i \Le P$, the $k$-nearest-neighbor radius of $x_i$ is
\begin{equation}
    \mathrm{NND}_k(x_i;\mathbf x) = \inf\Big\{ r \Ge 0 :
    \# \big\{1\Le j \Le P: 0 < |x_i - x_j| < r\big\} \Ge k
    \Big\},
\end{equation}
i.e., the distance from $x_i$ to its $k$-th nearest neighbor in $\mathbf x$. The coverage of $\mathbf x$ by the pushforward samples $\mathbf y$ is then
\begin{equation}
    \mathrm{Coverage}_k(\mathbf y;\mathbf x) =
    \frac{1}{P} \sum_{i=1}^P \mathbf 1\Bigl\{
        \exists\, 1\Le j\Le N : |y_j - x_i| \Le \mathrm{NND}_k(x_i;\mathbf x)
    \Bigr\},
\end{equation}
i.e., the fraction of points in $\mathbf x$ whose $k$-nearest-neighbor ball contains at least one pushforward sample. A coverage close to $1$ certifies that every region populated by $\hat\mu_1$ has been reached by the flow, while a coverage substantially below $1$ exposes precisely the missing modes that the ESS alone cannot see.

\bibliographystyle{plainnat}
\bibliography{references}

\end{document}